% (c) Tim Dokchitser, redlib library, v3.1, October 2024, https://people.maths.bris.ac.uk/~matyd/redlib/
%
% Example TeX file for using redlib.sty to display reduciton type labels

\documentclass[12pt]{amsart}
\usepackage{redlib}
\begin{document}

This is an example TeX file for the \texttt{redlib} library.
It uses redlib.sty, which relies on standard packages pbox, tikz and ams packages,
and implements \texttt{$\backslash$redtype\{...\}} macro and tikz defaults for special fibres.\\[5pt]

\begin{verbatim}
AttachSpec("redlib.spec");    // Magma
R:=ReductionType("11^5,8,9-{9-5}(3)9^1,5,3-{3-3}(1)T-{1-8}-{2-5}c1");
TeX(DualGraph(R));
\end{verbatim}

\begin{tikzpicture}[xscale=1,yscale=0.9]
\draw[thick,shorten >=2] (4/15,0)--(157/30,0) node[inner sep=2pt,blue,above left,scale=0.8] {11} node[inner sep=2pt,below left,scale=0.8] {$\Gamma_1$};
\draw[shorten <=-3pt, shorten >=-3pt] (3/5,0)--node[inner sep=2pt,scale=0.8,blue,right] {9} (3/5,3/5);
\draw[shorten <=-3pt, shorten >=-3pt] (3/5,3/5)--node[inner sep=2pt,scale=0.8,blue,above] {7} (0,3/5);
\draw[shorten <=-3pt, shorten >=-3pt] (0,3/5)--node[inner sep=2pt,scale=0.8,blue,left] {5} (0,6/5);
\draw[shorten <=-3pt, shorten >=-3pt] (0,6/5)--node[inner sep=2pt,scale=0.8,blue,above] {3} (3/5,6/5);
\draw[shorten <=-3pt, shorten >=-3pt] (3/5,6/5)--node[inner sep=2pt,scale=0.8,blue,right] {1} (3/5,9/5);
\draw[shorten <=-3pt, shorten >=-3pt] (3/5,9/5)--node[inner sep=2pt,scale=0.8,blue,above] {1} (0,9/5);
\draw[shorten <=-3pt, shorten >=-3pt] (0,9/5)--node[inner sep=2pt,scale=0.8,blue,left] {1} (0,12/5);
\draw[shorten <=-3pt, shorten >=-3pt] (0,12/5)--node[inner sep=2pt,scale=0.8,blue,above] {1} (3/5,12/5);
\draw[shorten <=-3pt, shorten >=-3pt] (3/5,12/5)--node[inner sep=2pt,scale=0.8,blue,right] {5} (3/5,3);
\draw[shorten <=-3pt, shorten >=-3pt] (2,0)--node[inner sep=2pt,scale=0.8,blue,right] {5} (2,36/49);
\draw[shorten <=-3pt, shorten >=-3pt] (2,36/49)--node[inner sep=2pt,scale=0.8,blue,above] {4} (7/5,36/49);
\draw[shorten <=-3pt, shorten >=-3pt] (7/5,36/49)--node[inner sep=2pt,scale=0.8,blue,left] {3} (7/5,72/49);
\draw[shorten <=-3pt, shorten >=-3pt] (7/5,72/49)--node[inner sep=2pt,scale=0.8,blue,below right=-1pt and -1pt] {2} (37/20,99/49);
\draw[shorten <=-3pt, shorten >=-3pt] (17/5,0)--node[inner sep=2pt,scale=0.8,blue,right] {8} (17/5,36/49);
\draw[shorten <=-3pt, shorten >=-3pt] (17/5,36/49)--node[inner sep=2pt,scale=0.8,blue,above] {5} (14/5,36/49);
\draw[shorten <=-3pt, shorten >=-3pt] (14/5,36/49)--node[inner sep=2pt,scale=0.8,blue,left] {2} (14/5,72/49);
\draw[shorten <=-3pt, shorten >=-3pt] (14/5,72/49)--node[inner sep=2pt,scale=0.8,blue,below right=-1pt and -1pt] {1} (13/4,99/49);
\draw[thick,shorten >=2] (91/60,99/49)--(61/12,99/49) node[inner sep=2pt,blue,above left,scale=0.8] {3} node[inner sep=2pt,below left,scale=0.8] {$\Gamma_3$};
\draw[shorten <=-3pt, shorten >=-3pt] (49/20,99/49)--node[inner sep=2pt,scale=0.8,blue,right] {3} (49/20,3);
\draw[thick,shorten >=2] (-1/3,3)--(257/60,3) node[inner sep=2pt,blue,above left,scale=0.8] {9} node[inner sep=2pt,below left,scale=0.8] {$\Gamma_2$};
\draw[shorten <=-3pt] (0,3)--node[inner sep=2pt,scale=0.8,blue,right] {1} (0,18/5);
\end{tikzpicture}

\begin{verbatim}
Label(R: tex);
\end{verbatim}

\redtype{11^{5,8,9}\e{9,5}(3)9^{1,5,3}\e{3,3}(1)T\e{1,8}\e{2,5}c_1}

\end{document}
